\documentclass{article}
\usepackage[margin=0.7in]{geometry}
\usepackage{amsmath,array,feynman-fun}
\pagestyle{empty}
\makeatletter
\newcount\perfassertions
\newcommand{\PerfAssert}[3]{%
  \corr@getlistitemorblank{#1}{#2}%
  \edef\perfexpected{#3}%
  \ifx\corr@currentlistitem\perfexpected\else
    \PackageError{performance-regression}{List lookup changed at slot #2}
      {Expected \meaning\perfexpected; got \meaning\corr@currentlistitem}%
  \fi
  \global\advance\perfassertions by 1
}
\newcommand{\PerfListChecks}{%
  \begingroup
  \corr@preparelists
  \def\PerfLeft{ a,,x_{1,2},d }%
  \def\PerfRight{i,j,k,l}%
  \PerfAssert{\PerfLeft}{1}{a}%
  \PerfAssert{\PerfRight}{1}{i}%
  \PerfAssert{\PerfLeft}{2}{}%
  \PerfAssert{\PerfRight}{2}{j}%
  \PerfAssert{\PerfLeft}{3}{x_{1,2}}%
  \PerfAssert{\PerfRight}{3}{k}%
  \PerfAssert{\PerfLeft}{4}{d}%
  \PerfAssert{\PerfLeft}{8}{}%
  \PerfAssert{}{1}{}%
  \def\PerfLeft{u,v,w}% % Same input macro, different contents.
  \PerfAssert{\PerfLeft}{2}{v}%
  \PerfAssert{r,s,t}{3}{t}% % Evict either previous parse.
  \PerfAssert{\PerfRight}{4}{l}%
  \begingroup
    \def\PerfRight{nested,scope}%
    \PerfAssert{\PerfRight}{2}{scope}%
  \endgroup
  \PerfAssert{\PerfRight}{2}{j}%
  \def\PerfValue{A}\def\PerfLeft{\PerfValue,B}%
  \PerfAssert{\PerfLeft}{1}{A}%
  \def\PerfValue{C}%
  \PerfAssert{\PerfLeft}{1}{C}%
  \ignoreemptyitems\corr@preparelists
  \PerfAssert{a,,b}{2}{b}%
  \reademptyitems\corr@preparelists
  \PerfAssert{a,,b}{2}{}%
  \PerfAssert{a,,b}{3}{b}%
  \endgroup
  \typeout{PASS performance compatibility: \the\perfassertions\space list assertions}%
}
\makeatother
\newcommand{\PerfBox}[1]{\BoxLoopCorr[
  scale=.68,show-momenta=false,
  box-propagator-left-indices={a,b,c,d,e,f,g,h},
  box-propagator-right-indices={i,j,k,l,m,n,o,p},#1]}
\begin{document}
\PerfListChecks
% Exercise style definitions whose bodies themselves accept an argument.
\begin{tikzpicture}\begin{feynhand}
\vertex (A) at (0,0);\vertex (B) at (1.5,0);\vertex (C) at (3,0);
\propag[every photon={/tikz/draw=pink},photon] (A) to (B);
\propag[every proca={/tikz/draw=blue},proca] (B) to (C);
\end{feynhand}\end{tikzpicture}
\begin{center}
\begin{tabular}{cc}
\PerfBox{box-line=plain} & \PerfBox{box-line={ewboson,circ=W}}\\[12pt]
\BoxLoopCorr[scale=.68,show-momenta=false] &
\FourPtCorr[box-top-line={proca,pink},topology=box,scale=.68,show-momenta=false]\\[12pt]
\FourPointCorr[scale=.68,topology=box,box-top-line={proca,pink},show-momenta=false] &
\CrossedVertexCorr[scale=.68,crossed-vertex-orientation=left,
 triangle-contact-internal-momenta={k,q-k,q},
 triangle-contact-propagator-left-indices={a,b,c,d,e,f,g},
 triangle-contact-propagator-right-indices={h,i,j,k,l,m,n}]
\end{tabular}
\end{center}
\newpage
\begingroup
\corrset{scalar-arrow-direction=out,color=purple,momentum-labels={k_1,k_2,k_3}}
\[\ThreePtCorr[scale=.72,line=spol]\qquad
  \ThreePointCorr[scale=.72,line=spol,scalar-arrow-direction=in]\]
\endgroup
\[\ThreePointCorr[scale=.72,line=spol]\qquad
  \ThreePointCorr[scale=.72,line=proca,momentum-labels={,,}]\]
\[\begin{array}{cc}
\TTreeChannelCorr[scale=.68,insertion=bubble,insertion-slot=2] &
\STreeChannelCorr[scale=.68,insertion=triangle]
\end{array}\]
\[\begin{aligned}
 &\FourPointCorr[scale=.68,insertion=seagull,insertion-slot=3]
 \qquad\FourPointCorr[scale=.68,show-momenta=false]
\end{aligned}\]
\end{document}
